\lyseis{Γ3.1}{Πειράματα τύχης και δειγματικός χώρος}
\sol{1} Πειράματα τύχης είναι τα α), γ), ε). Το β) έχει πάντα αποτέλεσμα $42$, το δ) το ξέρουμε (είναι η ώρα που δείχνει το ρολόι) και το στ) είναι σίγουρο ότι το νερό παγώνει στους $-5$ $^\circ$C.
\sol{2} α) $\{\mathrm{Κ},\ \mathrm{Γ}\}$. β) $\{1,2,3,4,5,6\}$. γ) $\{\text{άσπρη},\ \text{μαύρη},\ \text{κόκκινη}\}$. δ) $\{\text{άρτιος},\ \text{περιττός}\}$.
\sol{3} α) $\{\text{Τρίτη},\ \text{Τετάρτη}\}$, $2$ στοιχεία. β) $\{0,2,4\}$ (το $2$ εμφανίζεται δύο φορές αλλά μετράει ως ένα αποτέλεσμα), $3$ στοιχεία. γ) Τα διαφορετικά γράμματα είναι $\{\mathrm{Π},\ \mathrm{Ι},\ \mathrm{Θ},\ \mathrm{Α},\ \mathrm{Ν},\ \mathrm{Ο},\ \mathrm{Τ},\ \mathrm{Η}\}$, $8$ στοιχεία.
\sol{4} α) $\{6,\ \text{όχι }6\}$. β) Τα υπόλοιπα της διαίρεσης με το $3$ για τους αριθμούς $1,\dots,6$ είναι $1,2,0,1,2,0$, άρα $\{0,1,2\}$. γ) Πρώτοι είναι οι $2,3,5$. Δειγματικός χώρος $\{\text{πρώτος},\ \text{όχι πρώτος}\}$.
\sol{5} Παραδείγματα: α) ο τροχός του λούνα παρκ με οκτώ ίσα τμήματα αριθμημένα από $1$ ως $8$. β) Ρωτάω τυχαία έναν συμμαθητή αν του αρέσουν τα Μαθηματικά και σημειώνω ΝΑΙ ή ΟΧΙ. γ) Επιλέγω τυχαία μία από τις τρεις πρώτες μέρες της εβδομάδας. Υπάρχουν και άλλες απαντήσεις.
\sol{6} α) Λάθος. β) Σωστό. γ) Λάθος: είναι διαφορετικά αποτελέσματα (διαφορετική σειρά). δ) Σωστό. ε) Λάθος: γράφουμε όλα τα δυνατά αποτελέσματα.
\sol{7} α) $\{\mathrm{ΦΛ},\ \mathrm{ΛΦ}\}$. β) $\{\mathrm{ΦΦ},\ \mathrm{ΦΛ},\ \mathrm{ΛΦ},\ \mathrm{ΛΛ}\}$.
\sol{8} $\{\mathrm{ΝΝ},\ \mathrm{ΝΙ},\ \mathrm{ΝΗ},\ \mathrm{ΙΝ},\ \mathrm{ΙΙ},\ \mathrm{ΙΗ},\ \mathrm{ΗΝ},\ \mathrm{ΗΙ},\ \mathrm{ΗΗ}\}$, $9$ στοιχεία ($3$ επιλογές για τον πρώτο αγώνα και $3$ για τον δεύτερο).
\sol{9} $\{0,2,4,6,8\}$.
\sol{10} $\{33,37,73,77\}$.
\sol{11} $\{\mathrm{ΤΑ},\ \mathrm{ΤΚ},\ \mathrm{ΦΑ},\ \mathrm{ΦΚ}\}$.
\sol{12} $\{\mathrm{ΚΚΚ},\ \mathrm{ΚΚΓ},\ \mathrm{ΚΓΚ},\ \mathrm{ΚΓΓ},\ \mathrm{ΓΚΚ},\ \mathrm{ΓΚΓ},\ \mathrm{ΓΓΚ},\ \mathrm{ΓΓΓ}\}$, $8$ στοιχεία.
\sol{13} Όχι. Ο δειγματικός χώρος είναι ο ίδιος κάθε φορά: $\{\mathrm{Κ},\ \mathrm{Γ}\}$. Και τα δύο αποτελέσματα είναι δυνατά ξανά, ό,τι κι αν έγινε πριν. Το νόμισμα δεν θυμάται τις προηγούμενες ρίψεις.
\sol{14} Αν το πρώτο ζάρι φέρει $1$, το δεύτερο μπορεί να φέρει $1,2,3,4,5,6$: $6$ αποτελέσματα $(1,1),\dots,(1,6)$. Το ίδιο ισχύει για καθεμία από τις $6$ έδρες του πρώτου. Άρα $6\cdot6=36$ αποτελέσματα, όχι $12$.
\sol{15} Ναι. Π.χ. αν ρίξω ζάρι και σημειώσω «φέρνω αριθμό από $1$ έως $6$», ο δειγματικός χώρος είναι $\{\text{ναι}\}$. Τότε το αποτέλεσμα είναι σίγουρο και δεν υπάρχει τύχη, άρα ουσιαστικά δεν είναι πείραμα τύχης.
\sol{16} Ελεύθερη εργασία. Π.χ. «το πρώτο γράμμα του ονόματος του επόμενου που θα περάσει από την πόρτα» (πολλά στοιχεία), «βρέχει αύριο ή όχι» ($\{\text{ναι},\ \text{όχι}\}$), «το χρώμα του πρώτου αυτοκινήτου στον δρόμο».
\sol{17} $\{\text{κούπα},\ \text{καρό},\ \text{σπαθί},\ \text{μπαστούνι}\}$ ($4$ στοιχεία). Αν σημειώναμε ολόκληρο το φύλλο θα είχε $52$ στοιχεία.
